\chapter{L'autoregolazione come tratto transdiagnostico nella prima infanzia}
\label{cap:due}
    \section{Dalla regolazione tipica alla disregolazione}
    \label{sec:2_1}

        Se nel Capitolo 1 l'attenzione è stata posta sulle traiettorie tipiche di sviluppo dei processi di regolazione e sulla loro natura multidimensionale, nel presente capitolo tali processi verranno analizzati in relazione alle possibili difficoltà e deviazioni evolutive, con particolare riferimento al concetto di disregolazione e al suo ruolo come dimensione transdiagnostica nella prima infanzia.

        Alla luce delle traiettorie tipiche di sviluppo dei processi di regolazione descritte nel Capitolo 1, il funzionamento regolativo nella prima infanzia può essere compreso lungo un continuum evolutivo che va da modalità più adattive a forme progressivamente meno organizzate di modulazione delle emozioni, dell'attenzione e del comportamento. In questa prospettiva, sviluppo tipico e atipico non rappresentano condizioni qualitativamente distinte, ma esiti differenti di uno stesso processo evolutivo, che si differenziano per il grado di integrazione, flessibilità e stabilità dei sistemi regolativi nel tempo \cite{calkins2007,cicchetti2009}.

        La distinzione tra regolazione e disregolazione non si basa quindi sulla semplice presenza di specifici comportamenti, ma sul loro livello di organizzazione e sulla loro funzionalità rispetto alle richieste del contesto. In questa prospettiva, la disregolazione non si configura come la mera presenza di comportamenti problematici, ma come una difficoltà più ampia nella modulazione e nell'integrazione dei diversi sistemi coinvolti nel funzionamento regolativo, tra cui componenti emotive, attentive e comportamentali \cite{eisenberg2010}.

        Tali difficoltà possono manifestarsi con gradi diversi di stabilità e persistenza nel tempo e con differente sensibilità al contesto: mentre alcune espressioni risultano transitorie e coerenti con il livello di sviluppo, altre possono assumere un carattere più stabile, rigido e meno modulabile, configurandosi come possibili indicatori di rischio evolutivo. In questo senso, la valutazione del funzionamento regolativo richiede di considerare non solo la presenza dei comportamenti, ma anche la loro organizzazione, la loro adattività e la loro modificabilità nelle diverse situazioni \cite{tronick2007}.

        Tale approccio è coerente con le prospettive più recenti della psicopatologia dello sviluppo, che sottolineano la continuità tra funzionamento tipico e manifestazioni di difficoltà, in particolare nelle fasi precoci dello sviluppo, quando le competenze regolative risultano ancora emergenti e fortemente sensibili alle condizioni ambientali \cite{cicchetti2009}.

        In questa prospettiva, la disregolazione non rappresenta soltanto una deviazione dalle traiettorie tipiche di sviluppo, ma può configurarsi come un indicatore precoce di vulnerabilità evolutiva, rilevante per comprendere l'emergere di differenti esiti psicopatologici. In questo quadro, risulta fondamentale considerare la variabilità tipica che caratterizza i processi di regolazione nella prima infanzia, al fine di distinguere tra espressioni evolutivamente appropriate e possibili segnali di difficoltà, che verranno approfonditi nelle sezioni successive.

        \subsection{Variabilità nello sviluppo tipico}
        \label{sec:2_1_1}

            In continuità con quanto descritto nelle traiettorie tipiche di sviluppo (cfr. Capitolo 1, §1.6), la prima infanzia si caratterizza per un'elevata variabilità nel funzionamento regolativo, che rappresenta un elemento intrinseco dello sviluppo tipico. Tra i 18 e i 36 mesi, le competenze di regolazione sono ancora in fase di consolidamento e si manifestano in modo disomogeneo, con oscillazioni legate sia alle caratteristiche individuali sia alle condizioni contestuali.

            In questo periodo evolutivo è frequente osservare comportamenti quali crisi di rabbia, difficoltà nella tolleranza dell'attesa, oppositività o rapide fluttuazioni dello stato emotivo. Tali manifestazioni rientrano generalmente nella variabilità tipica e riflettono la natura ancora emergente dei processi di regolazione emotiva e comportamentale \cite{eisenberg2010,kochanska2001}.

            Un elemento caratteristico di questa fase riguarda la marcata sensibilità del comportamento alle condizioni ambientali. Il bambino può mostrare modalità di regolazione relativamente più organizzate in contesti familiari e prevedibili, ma incontrare difficoltà in situazioni nuove o emotivamente attivanti. Fattori quali la stanchezza, le transizioni tra attività o cambiamenti nelle routine quotidiane possono influenzare significativamente la capacità di modulare emozioni e comportamento, determinando fluttuazioni nel funzionamento regolativo (Tronick, 2007; Kochanska et al., 2001\textbf{).}

            Inoltre, nella prima infanzia, episodi di difficoltà nella regolazione possono essere seguiti da un ritorno relativamente rapido a uno stato di equilibrio, spesso facilitato dal supporto del caregiver. La possibilità di beneficiare dei processi di co-regolazione rappresenta un aspetto centrale del funzionamento tipico in questa fase dello sviluppo. In questa prospettiva, la responsività del bambino al supporto dell'adulto costituisce un indicatore rilevante della flessibilità del sistema regolativo, mentre una ridotta sensibilità alla co-regolazione può rappresentare un possibile segnale di difficoltà \cite{calkins2007,tronick2007}\textbf{.}

            La letteratura evidenzia come la comprensione del funzionamento regolativo nei primi anni di vita richieda un'attenzione specifica al contesto in cui i comportamenti si manifestano, nonché alla loro variabilità e modificabilità.

            Nel complesso, la variabilità tipica nei processi di regolazione riflette la natura dinamica dello sviluppo nella prima infanzia e rappresenta una base fondamentale per interpretare i comportamenti osservabili in modo evolutivamente appropriato \cite{calkins2007,kochanska2001}\textbf{.}

        \subsection{Manifestazioni della disregolazione nei primi anni di vita: caratteristiche quantitative, qualitative e contestuali}
        \label{sec:2_1_2}

            Alla luce della variabilità tipica descritta nella sezione precedente, risulta fondamentale delineare in modo più specifico le caratteristiche attraverso cui la disregolazione può manifestarsi nei primi anni di vita. Come già anticipato, la distinzione tra funzionamento regolativo tipico e segnali di difficoltà non si basa sulla semplice presenza di determinati comportamenti, ma su alcune dimensioni chiave che riguardano il modo in cui tali comportamenti si presentano, si organizzano e si modificano in relazione al contesto \cite{morris2007,tronick2011}.

            In linea con la letteratura sullo sviluppo della regolazione e sulla disregolazione precoce, è possibile distinguere tre principali livelli descrittivi utili per caratterizzare tali manifestazioni nella prima infanzia: dimensioni quantitative, qualitative e contestuali \cite{calkins2007,eisenberg2010}. In questa prospettiva, tali dimensioni assumono significato non come indicatori isolati, ma come espressione del livello di organizzazione del sistema regolativo nel suo complesso, inteso come insieme coordinato di processi emotivi, fisiologici e cognitivi \cite{beauchaine2019,beauchaine2015}.

            Dal punto di vista quantitativo, un primo elemento rilevante riguarda l'intensità delle risposte emotive e comportamentali. Nei bambini piccoli è frequente osservare reazioni emotive anche molto intense; tuttavia, nelle condizioni di possibile disregolazione tali risposte tendono a presentarsi con livelli di attivazione particolarmente elevati, difficilmente contenibili e sproporzionati rispetto alle richieste della situazione. In questo senso, l'intensità non rappresenta di per sé un indicatore di difficoltà, ma assume rilevanza quando riflette una ridotta capacità di regolazione dell'attivazione emotiva e fisiologica e una difficoltà nel recupero verso uno stato di equilibrio, configurandosi come espressione di un funzionamento regolativo meno stabile \cite{eisenberg2010,beauchaine2015}.

            Un secondo indicatore riguarda la frequenza degli episodi di difficoltà regolativa. Nello sviluppo tipico, manifestazioni di disregolazione possono emergere in modo occasionale e situazionale; al contrario, una frequenza elevata e ricorrente di tali episodi può indicare una difficoltà nella regolazione dinamica degli stati interni e una minore stabilità nel funzionamento regolativo nel tempo \cite{cole1994}. In questo senso, la frequenza non si limita a descrivere la quantità degli episodi, ma rappresenta un indice della continuità e coerenza del funzionamento regolativo.

            Un ulteriore elemento quantitativo è rappresentato dalla pervasività delle difficoltà, ovvero dalla loro presenza in diversi contesti e momenti della vita quotidiana. Quando le difficoltà regolative si manifestano in modo trasversale a situazioni differenti, indipendentemente dalle caratteristiche del contesto, esse possono riflettere una difficoltà più ampia nell'organizzazione del sistema regolativo, suggerendo una minore capacità di adattamento flessibile alle richieste ambientali \cite{cole1994}.

            Nel loro insieme, tali indicatori quantitativi assumono significato clinico non tanto come caratteristiche isolate, ma in quanto espressione della stabilità e del grado di coordinamento del funzionamento regolativo nel suo complesso \cite{calkins2007}.

            Accanto agli aspetti quantitativi, assumono particolare rilevanza anche le caratteristiche qualitative del funzionamento regolativo. In questo ambito, un elemento centrale riguarda il grado di flessibilità del comportamento. Nello sviluppo tipico, anche in presenza di difficoltà, il bambino mostra una certa variabilità nelle risposte e una capacità, seppur limitata, di modificare il proprio comportamento in funzione delle richieste dell'ambiente o del supporto dell'adulto. Al contrario, nelle manifestazioni di disregolazione si osserva spesso una maggiore rigidità, con risposte ripetitive, stereotipate o poco modulabili, difficilmente adattabili alle diverse situazioni \cite{morris2007}.

            Un ulteriore aspetto qualitativo riguarda la capacità di gestione degli stati emotivi e comportamentali. Nei quadri di maggiore difficoltà, le risposte appaiono meno differenziate e più globali, con una tendenza a mantenere elevati livelli di attivazione o a mostrare difficoltà nel recupero dell'equilibrio. Ciò riflette una limitata capacità di utilizzare strategie regolative efficaci e una ridotta integrazione tra i diversi sistemi coinvolti nel funzionamento regolativo \cite{tronick2011}.

            In questa prospettiva, la disregolazione non può essere interpretata esclusivamente come un aumento quantitativo di comportamenti problematici, ma come una alterazione qualitativa dell'organizzazione dei processi regolativi, che coinvolge il coordinamento tra sistemi emotivi, fisiologici e di controllo emergente \cite{beauchaine2019,blair2015}.

            Infine, un ulteriore livello di analisi riguarda le caratteristiche contestuali delle manifestazioni disregolative. Come evidenziato nella sezione precedente, nei primi anni di vita il comportamento del bambino è fortemente influenzato dal contesto e dal supporto dell'adulto. In condizioni di sviluppo tipico, le difficoltà regolative tendono a essere sensibili alle condizioni ambientali: il bambino può mostrare una maggiore organizzazione in contesti familiari e prevedibili e beneficiare della co-regolazione offerta dal caregiver \cite{tronick2011}.

            Nelle manifestazioni di disregolazione, invece, si osserva frequentemente una ridotta sensibilità alle richieste ambientali e al supporto dell'adulto. Le risposte del bambino possono risultare poco modulabili anche in presenza di contesti strutturati o di interventi regolativi esterni, suggerendo una difficoltà nel coordinamento tra processi interni e influenze ambientali e una limitata capacità di utilizzare il contesto come risorsa regolativa \cite{morris2007}.

            Queste dimensioni risultano particolarmente rilevanti anche sul piano clinico, in quanto consentono di individuare precocemente configurazioni di funzionamento potenzialmente associate a traiettorie di rischio evolutivo \cite{beauchaine2015,cole1994}. Tali dimensioni trovano inoltre riscontro nella letteratura empirica sulla disregolazione precoce, che evidenzia come elevata intensità, frequenza e scarsa modulabilità delle risposte emotive rappresentino indicatori rilevanti di vulnerabilità evolutiva.

            Nel complesso, le caratteristiche quantitative, qualitative e contestuali delle manifestazioni regolative forniscono quindi un quadro descrittivo utile per distinguere tra espressioni evolutivamente appropriate e possibili segnali di difficoltà nei primi anni di vita.

            Tale distinzione non implica una categorizzazione rigida tra normalità e patologia, ma si colloca all'interno di una prospettiva dimensionale che considera il funzionamento regolativo lungo un continuum di organizzazione e adattività. In questa prospettiva, la disregolazione può essere intesa non solo come espressione di una difficoltà evolutiva circoscritta, ma come una modalità di funzionamento che attraversa diversi domini dello sviluppo e che può contribuire a configurare differenti traiettorie di rischio, ponendo le basi per una sua interpretazione in chiave transdiagnostica \cite{beauchaine2019,astle2021}.

    \section{L'autoregolazione come tratto transdiagnostico nello sviluppo}
    \label{sec:2_2}

        Alla luce di quanto discusso nelle sezioni precedenti, la disregolazione nei primi anni di vita può essere interpretata non soltanto come espressione di possibili difficoltà evolutive, ma anche come una dimensione di funzionamento che attraversa diversi domini dello sviluppo e che può assumere rilevanza rispetto a molteplici esiti psicopatologici. In questa prospettiva, l'autoregolazione può essere considerata un costrutto centrale per comprendere la variabilità delle traiettorie evolutive, collocandosi al crocevia tra processi emotivi, cognitivi e comportamentali.

        Come discusso nel Capitolo 1, l'autoregolazione è stata concettualizzata come un costrutto multidimensionale che emerge dall'integrazione di componenti emotive, cognitive, comportamentali e fisiologiche \cite{calkins2007,blair2015}. Questa natura articolata ha portato a proporre l'autoregolazione come un costrutto rilevante in chiave transdiagnostica, in quanto coinvolge processi di base che attraversano diversi domini dello sviluppo e risultano implicati sia nel funzionamento tipico sia in quello atipico \cite{beauchaine2019}.

        Questo modo di concettualizzare il funzionamento regolativo si inserisce all'interno di un più ampio cambiamento di paradigma nella psicopatologia contemporanea, che ha progressivamente messo in discussione i modelli diagnostici categoriali tradizionali. In questa prospettiva, l'attenzione si è progressivamente spostata dalle categorie diagnostiche ai processi di base condivisi tra differenti condizioni, favorendo una lettura del funzionamento psicologico in termini dimensionali piuttosto che categoriali \cite{astle2021}.

        Applicata allo sviluppo precoce, questa prospettiva assume una rilevanza specifica. Nei primi anni di vita, infatti, le competenze regolative sono ancora in fase di organizzazione e le manifestazioni comportamentali risultano spesso poco differenziate, rendendo meno adeguata una lettura in termini di categorie diagnostiche stabili \cite{calkins2007}. In questo contesto, l'analisi dei processi sottostanti, quali la capacità di modulare emozioni, attenzione e comportamento, consente di individuare dimensioni di funzionamento che attraversano differenti traiettorie evolutive e contribuiscono alla comprensione sia dello sviluppo tipico sia di quello atipico \cite{blair2015}.

        Considerare l'autoregolazione in chiave transdiagnostica non implica una sua sovrapposizione con specifiche condizioni patologiche, ma il riconoscimento del suo ruolo come dimensione fondamentale dell'organizzazione del comportamento, la cui variabilità può contribuire a spiegare la diversificazione degli esiti evolutivi. In questa prospettiva, la disregolazione può essere interpretata come una modalità di funzionamento che riflette una difficoltà nell'integrazione dei sistemi coinvolti nella regolazione e che, proprio per la sua natura trasversale, può assumere significato in relazione a molteplici traiettorie evolutive.

        In questa prospettiva, i processi di autoregolazione sono stati interpretati come meccanismi di base che mediano l'interazione tra sistemi neurobiologici e comportamentali, contribuendo alla vulnerabilità transdiagnostica \cite{beauchaine2019}.

        Va tuttavia sottolineato che l'approccio transdiagnostico non sostituisce completamente i modelli categoriali, che continuano a svolgere una funzione rilevante in ambito clinico e diagnostico, ma si propone piuttosto come prospettiva complementare per comprendere i meccanismi sottostanti allo sviluppo atipico.

        Questa prospettiva è stata ulteriormente sviluppata in diversi modelli dimensionali della psicopatologia contemporanea (es. modelli gerarchici e fattoriali), che verranno approfonditi nella sezione 2.4, i quali hanno evidenziato l'esistenza di dimensioni di funzionamento condivise tra differenti quadri clinici \cite{caspi2014,kotov2017}.

        Alla luce di queste considerazioni, risulta necessario definire in modo più preciso cosa si intenda per tratto transdiagnostico e quali siano le caratteristiche che ne delimitano il significato all'interno della psicopatologia dello sviluppo. Prima di applicare tale prospettiva al costrutto di autoregolazione, è quindi opportuno chiarire i criteri teorici che consentono di distinguere i processi propriamente transdiagnostici da altre dimensioni di funzionamento solo apparentemente condivise tra differenti quadri clinici. Questo aspetto verrà approfondito nella sezione seguente.

        \subsection{Definizione e caratteristiche dei tratti transdiagnostici}
        \label{sec:2_2_1}

            Sebbene il concetto di transdiagnosticità sia ampiamente utilizzato nella letteratura contemporanea, esso non è privo di ambiguità teoriche e richiede una definizione più precisa, al fine di evitare un uso eccessivamente estensivo o meramente descrittivo.

            In termini generali, la transdiagnosticità si riferisce alla presenza di processi, meccanismi o dimensioni di funzionamento che risultano comuni a più disturbi e che contribuiscono, in misura variabile, alla loro espressione \cite{astle2021}. Tuttavia, non tutti gli elementi condivisi tra differenti quadri psicopatologici possono essere considerati propriamente transdiagnostici.

            A questo proposito, è necessario distinguere tra livelli concettualmente differenti. I processi transdiagnostici fanno riferimento a meccanismi implicati direttamente nell'organizzazione del funzionamento psicopatologico e che operano trasversalmente a più condizioni cliniche \cite{nolenhoeksema2011}. Diversamente, i fattori di rischio rappresentano condizioni che aumentano la probabilità di esiti disadattivi, senza costituire necessariamente un meccanismo esplicativo comune. Infine, i correlati non specifici possono essere osservati in una pluralità di quadri clinici, ma non svolgono un ruolo funzionale nei processi sottostanti e non contribuiscono direttamente alla loro organizzazione.

            In questa prospettiva, un tratto può essere definito propriamente transdiagnostico quando non si limita a comparire in più condizioni cliniche, ma risulta implicato nei meccanismi sottostanti a differenti quadri psicopatologici e contribuisce in modo sistematico all'organizzazione del funzionamento psicologico\textbf{.} Questa concezione è coerente con la letteratura sui processi transdiagnostici, che sottolinea la necessità di identificare meccanismi funzionali condivisi tra differenti condizioni, piuttosto che semplici correlati descrittivi \cite{nolenhoeksema2011}. Ciò implica che tale tratto partecipi, almeno in parte, ai processi causali o di mantenimento della psicopatologia, piuttosto che rappresentare un semplice correlato descrittivo. In questo senso, la sola presenza trasversale non è sufficiente a definirne la natura transdiagnostica, ma è necessario che il tratto operi come dimensione funzionale all'interno dei processi che sostengono il funzionamento psicologico disadattivo.

            Le caratteristiche dei tratti transdiagnostici possono essere ulteriormente comprese considerando diversi livelli di analisi. A un primo livello, di natura descrittiva, essi si caratterizzano per la loro trasversalità, risultando osservabili in una pluralità di condizioni cliniche. Tuttavia, tale trasversalità non è di per sé sufficiente a definirne il significato, poiché uno stesso tratto può assumere configurazioni differenti e svolgere funzioni diverse nei vari quadri psicopatologici \cite{astle2021}. A un livello più propriamente teorico, i tratti transdiagnostici si distinguono per la loro non specificità diagnostica: essi non identificano un disturbo in senso stretto, ma rappresentano dimensioni di funzionamento che attraversano categorie differenti e contribuiscono alla loro organizzazione. Questo aspetto risulta particolarmente rilevante nella prima infanzia, fase in cui le manifestazioni del funzionamento psicologico appaiono ancora poco differenziate e difficilmente riconducibili a configurazioni diagnostiche stabili \cite{calkins2007}.

            Infine, sul piano funzionale, tali dimensioni riflettono processi implicati nell'adattamento dell'individuo alle richieste dell'ambiente e risultano strettamente connesse all'organizzazione del comportamento. In questa prospettiva, i tratti transdiagnostici possono essere intesi come dimensioni latenti che emergono dalla variabilità condivisa tra diverse manifestazioni del funzionamento cognitivo ed emotivo, collocandosi lungo un continuum che va da modalità più adattive a configurazioni progressivamente meno organizzate \cite{astle2021}.

            Ad esempio, difficoltà nei processi di controllo inibitorio possono essere osservate in diversi quadri del neurosviluppo e nei disturbi emotivo-comportamentali, assumendo configurazioni differenti in relazione all'organizzazione del funzionamento regolativo: nel disturbo da deficit di attenzione/iperattività (ADHD) tali difficoltà si esprimono prevalentemente come una ridotta capacità di inibire risposte impulsive e di modulare il comportamento in funzione delle richieste del contesto, mentre nei disturbi d'ansia possono manifestarsi come forme di ipercontrollo e rigidità comportamentale, associate a tendenze di evitamento e ritiro.

            Nonostante la sua utilità, il concetto di transdiagnosticità presenta alcune criticità. Il rischio di eccessiva generalizzazione può portare a includere sotto questa etichetta processi eterogenei nei loro meccanismi sottostanti, mentre la natura spesso latente di tali dimensioni pone rilevanti questioni metodologiche in termini di operazionalizzazione e misurazione empirica \cite{nolenhoeksema2011}. In particolare, la natura latente di tali dimensioni rende complessa la loro operazionalizzazione empirica, poiché differenti strumenti di misura possono catturare aspetti parzialmente sovrapponibili ma non equivalenti del funzionamento regolativo, sollevando interrogativi circa la validità costruttiva e la comparabilità dei risultati tra studi.

            In questo quadro, l'autoregolazione rappresenta una dimensione di particolare interesse. La sua natura multidimensionale e il suo coinvolgimento nell'integrazione di processi emotivi, cognitivi e comportamentali la rendono un candidato teoricamente rilevante per essere considerata in chiave transdiagnostica \cite{beauchaine2015}. Tuttavia, tale attribuzione richiede di essere valutata alla luce dei criteri sopra delineati, verificando in che misura i processi regolativi partecipino effettivamente ai meccanismi sottostanti le diverse traiettorie psicopatologiche.

            Per comprendere più a fondo questo aspetto, è necessario analizzare come il concetto di transdiagnosticità si sia sviluppato all'interno della psicopatologia dello sviluppo e quali trasformazioni teoriche ne abbiano favorito l'applicazione allo studio dei processi regolativi nella prima infanzia, tema che verrà approfondito nella sezione seguente.

        \subsection{Origine e sviluppo del concetto nella psicopatologia dello sviluppo}
        \label{sec:2_2_2}

            Negli ultimi decenni, il crescente interesse per l'approccio transdiagnostico si è sviluppato all'interno di una più ampia evoluzione della psicopatologia dello sviluppo, che ha progressivamente orientato l'attenzione verso l'analisi dei processi sottostanti alle traiettorie evolutive.

            Più che rappresentare una rottura netta rispetto ai modelli precedenti, tale cambiamento può essere compreso come il risultato di una trasformazione progressiva nel modo in cui vengono concettualizzati i meccanismi di funzionamento psicologico. In una fase iniziale, i modelli categoriali tradizionali hanno privilegiato l'identificazione di quadri clinici distinti e relativamente indipendenti, all'interno dei quali i processi di regolazione venivano considerati principalmente come competenze rilevanti per l'adattamento, ma non necessariamente interpretate come meccanismi condivisi tra differenti condizioni psicopatologiche \cite{beauchaine2019}.

            Un passaggio teorico fondamentale si colloca con l'affermarsi della developmental psychopathology, che concepisce lo sviluppo come il risultato di interazioni dinamiche tra fattori biologici, psicologici e ambientali e sottolinea la continuità tra funzionamento tipico e atipico \cite{cicchetti2006}. In questo quadro, l'attenzione si è progressivamente orientata verso dimensioni di funzionamento trasversali, capaci di spiegare la variabilità individuale lungo le traiettorie evolutive. Questo passaggio risulta particolarmente rilevante per la prospettiva transdiagnostica, poiché introduce l'idea che i processi psicologici possano essere compresi non in relazione a specifiche categorie diagnostiche, ma come dimensioni che attraversano differenti esiti evolutivi.

            In questa prospettiva, l'interesse per dimensioni di funzionamento trasversali riflette anche l'esigenza di individuare costrutti che risultino maggiormente informativi rispetto alle difficoltà funzionali quotidiane, superando i limiti descrittivi delle categorie diagnostiche tradizionali. I processi di regolazione hanno così iniziato a essere interpretati non solo come competenze adattive, ma come dimensioni centrali dell'organizzazione del comportamento, potenzialmente implicate in molteplici traiettorie di sviluppo.

            Parallelamente, i contributi provenienti dalla psicologia cognitiva e dalle neuroscienze dello sviluppo hanno consentito una progressiva specificazione dei meccanismi sottostanti ai processi regolativi, in particolare attraverso lo studio delle funzioni esecutive e dei sistemi di controllo del comportamento. L'integrazione tra approcci relazionali e neurocognitivi ha favorito una lettura multilivello del funzionamento regolativo, inteso come risultato dell'interazione tra sistemi emotivi, cognitivi e motivazionali \cite{blair2015,diamond2013}. In relazione alla prospettiva transdiagnostica, tale integrazione assume un ruolo centrale, in quanto permette di concepire l'autoregolazione come un processo di base che opera a diversi livelli di analisi e che contribuisce all'organizzazione del funzionamento psicologico in modo trasversale rispetto alle diverse condizioni cliniche.

            In questa direzione, difficoltà nei processi di autoregolazione, in particolare nei sistemi di controllo esecutivo, sono state progressivamente considerate come indicatori transdiagnostici di sviluppo atipico, in quanto implicate in molteplici traiettorie psicopatologiche \cite{beauchaine2015}. Ciò ha contribuito a spostare l'attenzione da una concezione dell'autoregolazione come competenza specifica, spesso considerata in relazione a singoli domini dello sviluppo, a una sua interpretazione come processo di base implicato nell'organizzazione generale del comportamento.

            Il successivo sviluppo di prospettive dimensionali nello studio della psicopatologia ha ulteriormente rafforzato questa direzione come evidenziato da recenti modelli gerarchici della psicopatologia che hanno messo in luce l'esistenza di dimensioni di funzionamento condivise tra differenti quadri clinici \cite{caspi2014,kotov2017}. In questo senso, l'interesse per i meccanismi sottostanti, piuttosto che per le categorie diagnostiche, ha favorito l'emergere di modelli orientati all'identificazione di dimensioni di funzionamento comuni \cite{astle2021}.

            All'interno di questo quadro, l'autoregolazione emerge come uno dei costrutti più rilevanti per l'analisi dei processi transdiagnostici, in quanto coinvolge meccanismi di base implicati nella modulazione delle emozioni, dell'attenzione e del comportamento e contribuisce alla comprensione della variabilità degli esiti evolutivi \cite{calkins2007}.

            Tuttavia, questa evoluzione teorica non è esente da criticità. Il rischio di eccessiva genericità e la difficoltà di definire con precisione i meccanismi coinvolti pongono importanti sfide concettuali e metodologiche \cite{nolenhoeksema2011}. In particolare, risulta necessario distinguere tra dimensioni effettivamente implicate nei processi psicopatologici e caratteristiche solo apparentemente condivise, evitando di attribuire un valore esplicativo a fenomeni che potrebbero riflettere semplici sovrapposizioni descrittive.

            In sintesi, l'evoluzione del concetto può essere letta come un passaggio progressivo: da una concezione della regolazione come competenza specifica, a una sua interpretazione come processo di base, fino alla sua attuale considerazione come possibile dimensione transdiagnostica. Questo percorso teorico consente di comprendere come l'autoregolazione sia progressivamente divenuta un costrutto centrale per l'analisi dei meccanismi condivisi nella psicopatologia dello sviluppo.

            Le evidenze relative al ruolo dei processi regolativi nei diversi ambiti del funzionamento psicopatologico verranno approfondite nelle sezioni successive.

    \section{Evidenze empiriche del ruolo transdiagnostico della disregolazione}
    \label{sec:2_3}

        In linea con quanto discusso nella sezione precedente, i modelli categoriali tradizionali mostrano limiti nel cogliere la complessità e l'eterogeneità dei profili di funzionamento nei bambini, favorendo l'emergere di prospettive dimensionali e transdiagnostiche \cite{cicchetti2009,astle2022}. In questa prospettiva, le difficoltà evolutive tendono a organizzarsi lungo dimensioni continue piuttosto che all'interno di categorie discrete, risultando frequentemente associate a pattern di funzionamento che attraversano diversi quadri diagnostici.

        Tra i costrutti che hanno ricevuto crescente attenzione in questo ambito, la disregolazione si configura come una dimensione particolarmente rilevante per comprendere il funzionamento tipico e atipico nei primi anni di vita. Come discusso nel Capitolo 1, i processi di regolazione emergono dall'organizzazione progressiva di sistemi cognitivi, emotivi, fisiologici e comportamentali, che contribuiscono congiuntamente alla modulazione dell'esperienza e dell'azione in relazione alle richieste del contesto \cite{calkins2007,blair2015}.

        In questa prospettiva, la disregolazione può essere intesa come una difficoltà nel coordinamento tra sistemi di risposta reattivi e processi di controllo volontario, con possibili ricadute sull'adattamento del bambino alle richieste ambientali. Numerosi contributi teorici ed empirici sottolineano come tali difficoltà non si presentino in modo isolato, ma tendano a coinvolgere più domini di funzionamento, contribuendo alla co-occorrenza di difficoltà cognitive, emotive e comportamentali \cite{beauchaine2019,astle2022}.

        Ulteriori evidenze indicano come i processi di regolazione rappresentino un nodo centrale nell'organizzazione dello sviluppo, risultando strettamente connessi sia ai sistemi emotivo-motivazionali sia ai processi di controllo cognitivo, tra cui le funzioni esecutive \cite{diamond2013,zelazo2020}. In questa prospettiva, la disregolazione può essere interpretata come espressione di una minore efficacia nei meccanismi che sostengono l'adattamento flessibile al contesto, piuttosto che come una caratteristica specifica di singoli disturbi.

        Coerentemente con queste prospettive, approcci recenti nell'ambito della psicopatologia dello sviluppo hanno proposto modelli transdiagnostici che enfatizzano il ruolo di dimensioni di funzionamento condivise tra diversi disturbi. In questa linea si colloca, ad esempio, il modello del~\emph{p factor}, che suggerisce l'esistenza di una dimensione generale di vulnerabilità psicopatologica associata a difficoltà pervasive nei processi di regolazione e controllo \cite{caspi2014}.

        Nel complesso, tali evidenze supportano una lettura della disregolazione come dimensione trasversale rilevante per la comprensione delle traiettorie di sviluppo tipico e atipico. In questa prospettiva, l'analisi dei diversi contesti evolutivi in cui emergono difficoltà regolative consente di approfondire le modalità attraverso cui tali processi si manifestano e si organizzano nel corso dello sviluppo.

        Alla luce di queste considerazioni, nelle sezioni seguenti verrà approfondito il ruolo della disregolazione in specifici ambiti dello sviluppo, a partire dai disturbi del neurosviluppo.

        \subsection{Disregolazione nei disturbi del neurosviluppo}
        \label{sec:2_3_1}

            Alla luce delle evidenze presentate, i disturbi del neurosviluppo rappresentano un ambito particolarmente rilevante per approfondire il ruolo della disregolazione nei primi anni di vita. Tali condizioni si caratterizzano infatti per un'elevata eterogeneità dei profili di funzionamento e per una frequente sovrapposizione tra difficoltà cognitive, emotive e comportamentali, rendendo particolarmente evidente il contributo dei processi di regolazione nell'adattamento del bambino al contesto \cite{astle2022,blair2015}.

            In questo ambito, difficoltà nella modulazione dell'attenzione, delle emozioni e del comportamento risultano frequentemente osservabili, configurandosi come espressione di una minore coordinazione tra processi reattivi e sistemi di controllo volontario, in linea con quanto descritto nello sviluppo dei processi di autoregolazione \cite{blair2015,diamond2013}. Tale configurazione risulta coerente con una lettura transdiagnostica delle difficoltà evolutive, in cui la disregolazione emerge come dimensione trasversale che attraversa diversi quadri clinici. Tra le condizioni maggiormente studiate in relazione ai processi di disregolazione rientrano il disturbo da deficit di attenzione/iperattività (ADHD), i disturbi dello spettro autistico e i disturbi del linguaggio.

            Nel disturbo da deficit di attenzione/iperattività (ADHD), le difficoltà di regolazione si manifestano in modo particolarmente evidente nei domini attentivo e comportamentale, ma coinvolgono anche la modulazione degli stati emotivi. In questo quadro, i comportamenti di disattenzione, impulsività e iperattività possono essere interpretati non soltanto come sintomi isolati, ma come espressione di una difficoltà più ampia nella modulazione delle risposte in funzione delle richieste del contesto.

            Numerose evidenze indicano che i bambini con ADHD presentano difficoltà nei processi di controllo inibitorio e nella regolazione dell'attenzione, che risultano fondamentali per sostenere comportamenti orientati allo scopo e per modulare le risposte impulsive \cite{diamond2013,zelazo2020,barkley1997}. In questa prospettiva, il modello di Barkley \cite{barkley1997} evidenzia come il deficit nei meccanismi inibitori rappresenti un elemento centrale nel disturbo, influenzando a cascata l'organizzazione del comportamento, la pianificazione dell'azione e la capacità di mantenere attivi gli obiettivi nel tempo.

            A questa interpretazione si affianca il dual pathway model proposto da Sonuga-Barke \cite{sonugabarke2005}, secondo cui le difficoltà osservate nell'ADHD possono derivare da alterazioni in sistemi parzialmente distinti, uno legato al controllo esecutivo e uno ai processi motivazionali. In particolare, la difficoltà nel modulare la sensibilità alla ricompensa e alla gratificazione immediata può contribuire a una minore capacità di mantenere comportamenti orientati a obiettivi a lungo termine, aumentando la variabilità delle risposte in funzione del contesto.

            In entrambi i modelli, emerge il ruolo della disregolazione come difficoltà nel coordinamento tra processi di controllo top-down e risposte più automatiche e reattive, con una conseguente minore flessibilità nell'adattare il comportamento alle richieste ambientali. Accanto a questi aspetti, diversi studi evidenziano come nei bambini con ADHD siano frequentemente presenti difficoltà nella regolazione emotiva, che si manifestano, ad esempio, in una maggiore reattività alla frustrazione e in una ridotta capacità di modulare l'intensità e la durata delle risposte emotive \cite{eisenberg2010,beauchaine2019}. Tali difficoltà risultano particolarmente evidenti nelle situazioni che richiedono attesa, tolleranza della frustrazione o adattamento a richieste esterne, suggerendo una minore integrazione tra sistemi cognitivi e processi emotivo-motivazionali.

            Nel complesso, il profilo osservabile nell'ADHD può essere compreso come espressione di una disregolazione nei sistemi che coordinano attenzione, emozione e comportamento, piuttosto che come insieme di manifestazioni indipendenti. Questo quadro appare coerente con una prospettiva che interpreta l'ADHD non solo come un disturbo specifico, ma anche come una configurazione particolare di difficoltà nei processi di regolazione che possono essere rintracciati, con modalità diverse, in altri quadri evolutivi \cite{blair2015,zelazo2020}.

            In modo complementare, anche nei disturbi dello spettro autistico le difficoltà di regolazione si manifestano in maniera articolata e fortemente eterogenea, coinvolgendo in particolare la modulazione degli stati emotivi, la reattività agli stimoli e la flessibilità del comportamento. In questo contesto, la disregolazione può emergere sia come difficoltà nel modulare l'intensità e la durata delle risposte emotive, sia come alterazioni nei livelli di attivazione, con profili caratterizzati da iper-reattività o ipo-reattività agli stimoli ambientali.

            Numerose evidenze indicano che, in molti bambini con disturbo dello spettro autistico, la regolazione emotiva risulta meno flessibile e maggiormente dipendente dalle caratteristiche del contesto \cite{eisenberg2010,beauchaine2019,mazefsky2013}. Tali difficoltà si inseriscono all'interno di profili caratterizzati da una marcata variabilità individuale, coerentemente con modelli che sottolineano l'eterogeneità dei disturbi del neurosviluppo \cite{astle2022}.

            Un elemento particolarmente rilevante riguarda la modulazione sensoriale, frequentemente atipica in questo quadro, che contribuisce a influenzare i livelli di attivazione e la regolazione degli stati interni \cite{beauchaine2019,mazefsky2013}. La difficoltà nel filtrare o modulare gli input sensoriali può infatti amplificare o attenuare le risposte emotive e comportamentali, rendendo più complessa la regolazione degli stati di attivazione in funzione delle richieste ambientali.

            Parallelamente, la regolazione cognitiva appare coinvolta nella flessibilità comportamentale e nell'adattamento alle richieste del contesto, riflettendo una minore interazione funzionale tra sistemi di controllo e processi più automatici \cite{diamond2013,zelazo2020,samson2014}. In particolare, la difficoltà nel modificare strategie comportamentali o nel passare da un'attività all'altra può essere interpretata come espressione di una ridotta integrazione tra processi esecutivi e risposte reattive.

            La compresenza di questi fattori contribuisce a delineare profili in cui la disregolazione emerge dall'interazione tra reattività emotiva, processamento sensoriale e ridotta flessibilità cognitiva, configurandosi come un fenomeno multidimensionale. Questo suggerisce come le difficoltà osservate nei disturbi dello spettro autistico possano essere comprese non soltanto in termini di specificità del quadro clinico, ma anche come espressione di modalità particolari di organizzazione dei processi di regolazione, rilevanti in una prospettiva transdiagnostica.

            Una prospettiva parzialmente diversa emerge nei disturbi del linguaggio, in cui le difficoltà di regolazione risultano strettamente connesse al ruolo che il linguaggio svolge nello sviluppo dei processi di autoregolazione. Tali disturbi, definiti come disturbo primario del linguaggio, si caratterizzano per difficoltà persistenti nello sviluppo delle competenze linguistiche non attribuibili ad altre condizioni neurologiche o sensoriali \cite{bishop2017}.

            Numerosi contributi teorici ed empirici evidenziano come il linguaggio rappresenti uno strumento fondamentale per l'organizzazione del comportamento e la modulazione degli stati interni, in particolare attraverso il linguaggio interno e le forme di autoistruzione \cite{vygotsky1962,zelazo2015,vallotton2011}. In questa prospettiva, il linguaggio consente al bambino di guidare il comportamento in modo volontario, sostenere l'attenzione, rappresentare mentalmente le situazioni e mantenere attivi gli obiettivi nel tempo.

            Quando tali competenze risultano compromesse, può emergere una difficoltà nell'utilizzare strategie verbali per sostenere l'attenzione, pianificare l'azione e modulare le risposte comportamentali in funzione delle richieste del contesto. In particolare, una minore disponibilità di risorse linguistiche può incidere sulla capacità di utilizzare il linguaggio come strumento di mediazione del comportamento, riducendo la possibilità di organizzare e monitorare le proprie azioni in modo flessibile.

            Questo aspetto appare particolarmente rilevante nei contesti che richiedono controllo volontario, pianificazione o adattamento a richieste esterne, in cui il linguaggio svolge un ruolo chiave nel sostenere i processi di regolazione. Accanto a questi aspetti, il linguaggio svolge un ruolo rilevante anche nella regolazione emotiva, consentendo al bambino di comprendere, etichettare e modulare le proprie esperienze affettive \cite{eisenberg2010}. Una difficoltà in questo ambito può quindi riflettersi in una minore capacità di riconoscere e gestire gli stati emotivi, con possibili ricadute sull'adattamento al contesto.

            Nel complesso, i profili osservati nei disturbi del linguaggio possono essere interpretati come espressione di una disregolazione che coinvolge il coordinamento tra sistemi linguistici, cognitivi ed emotivi. In questa prospettiva, il ruolo del linguaggio come mediatore dei processi di regolazione consente di comprendere tali difficoltà come parte di un funzionamento più ampio, rilevante anche in altri quadri evolutivi \cite{blair2015,zelazo2015}.

            Nel complesso, l'analisi dei disturbi del neurosviluppo evidenzia come la disregolazione rappresenti una dimensione trasversale che coinvolge, seppur con modalità differenti, il coordinamento tra processi cognitivi, emotivi e comportamentali \cite{astle2022,beauchaine2019}. Le difficoltà di regolazione si configurano come espressione di una minore efficacia nei sistemi di controllo che sostengono l'adattamento alle richieste del contesto \cite{blair2015,zelazo2020}.

            Tale variabilità riflette la natura multidimensionale dei processi di autoregolazione e la loro progressiva organizzazione nel corso dello sviluppo \cite{calkins2007,diamond2013}. Allo stesso tempo, la presenza di difficoltà regolative in condizioni eterogenee suggerisce l'esistenza di meccanismi sottostanti comuni, rilevanti per la comprensione delle traiettorie di sviluppo atipico. In questa prospettiva, la disregolazione può essere considerata un elemento chiave per comprendere sia le specificità dei singoli quadri sia le somiglianze che attraversano diversi contesti evolutivi.

            Alla luce di queste considerazioni, risulta rilevante estendere l'analisi anche a condizioni di rischio evolutivo, come verrà approfondito nella sezione successiva.

        \subsection{Disregolazione in condizioni di rischio evolutivo e sviluppo atipico}
        \label{sec:2_3_2}

            Nel quadro delineato nelle sezioni precedenti, l'analisi della disregolazione nella prima infanzia assume particolare rilevanza nei contesti di rischio evolutivo e nelle condizioni di sviluppo atipico. In tali contesti, le difficoltà regolative tendono a manifestarsi con maggiore stabilità e pervasività, riflettendo specifiche configurazioni dello sviluppo neurobiologico e delle esperienze precoci, che incidono sull'organizzazione e sull'integrazione dei sistemi regolativi nel tempo \cite{cicchetti2009,calkins2007}.

            Tra le condizioni maggiormente studiate, il ritardo globale dello sviluppo e la nascita pretermine rappresentano due contesti particolarmente significativi per comprendere come i processi di regolazione possano risultare meno flessibili o meno coordinati nei primi anni di vita, pur all'interno di traiettorie evolutive tra loro differenti.

            Il ritardo globale dello sviluppo, che comprende condizioni come la sindrome di Down, è caratterizzato da un rallentamento generalizzato nei principali domini evolutivi, con implicazioni rilevanti anche per l'organizzazione dei processi di regolazione. In questi bambini, le difficoltà regolative emergono dall'interazione tra limitazioni nelle risorse cognitive, tempi di maturazione più lenti e caratteristiche specifiche del funzionamento neuropsicologico \cite{fidler2005,daunhauer2011}.

            La letteratura evidenzia un profilo regolativo disomogeneo, in cui competenze relativamente preservate sul piano socio-relazionale, quali la responsività all'interazione, coesistono con difficoltà nella modulazione attentiva e nel controllo comportamentale, in particolare nelle situazioni che richiedono il mantenimento dell'attenzione, la gestione della frustrazione o l'inibizione della risposta \cite{fidler2009}.

            Tali caratteristiche possono essere comprese, almeno in parte, alla luce di una minore efficienza dei processi di controllo esecutivo, in interazione con la più ampia organizzazione dei sistemi regolativi \cite{zelazo2020,lanfranchi2010}. I meccanismi di controllo top-down, deputati alla modulazione delle risposte emotive e comportamentali in funzione delle richieste contestuali, risultano infatti meno efficaci, rendendo più complessa l'organizzazione flessibile del comportamento. Ne deriva un funzionamento regolativo maggiormente sostenuto da processi bottom-up e dal supporto esterno, piuttosto che da forme autonome di controllo.

            Sul piano osservativo, tali difficoltà si esprimono in una minore capacità di sostenere e modulare l'attenzione e nella ridotta flessibilità comportamentale, con una maggiore dipendenza dalle strategie di co-regolazione offerte dall'adulto \cite{calkins2007}. Al contempo, il funzionamento regolativo risulta fortemente sensibile al contesto relazionale: interazioni prevedibili, strutturate e responsivamente sintonizzate possono infatti sostenere l'organizzazione del comportamento e favorire l'emergere di processi di regolazione più adattivi \cite{tronick2007}. Le difficoltà osservate vanno pertanto comprese all'interno di traiettorie di sviluppo specifiche, caratterizzate da tempi e modalità di integrazione differenti. In tale prospettiva, è inoltre importante considerare la significativa eterogeneità nei profili di funzionamento regolativo nei bambini ritardo globale dello sviluppo, con profili differenziati in funzione dell'interazione tra risorse individuali e contesti di sviluppo.

            Un ulteriore contesto di rischio evolutivo rilevante è rappresentato dalla nascita pretermine, che comporta un'interruzione anticipata dei processi di maturazione intrauterina e una conseguente esposizione precoce a un ambiente extrauterino spesso caratterizzato da elevata stimolazione e minore prevedibilità.

            Numerose evidenze indicano che i bambini nati pretermine presentano, già nei primi anni di vita, difficoltà nei processi di regolazione che coinvolgono domini emotivi, attentivi e comportamentali \cite{montagna2016}. Tali difficoltà riflettono l'interazione tra fattori neurobiologici, legati alla maturazione dei sistemi cerebrali implicati nella regolazione, ed esperienze precoci spesso caratterizzate da elevata stimolazione e discontinuità nelle interazioni con il caregiver, riflettendo una minore integrazione e coordinazione dei sistemi regolativi nei primi anni di vita.

            Dal punto di vista neuroevolutivo, la prematurità è associata a una minore integrazione dei circuiti che coordinano processi bottom-up e top-down, con conseguenze sulla capacità di modulare l'attivazione fisiologica ed emotiva \cite{blair2015}. Ciò si traduce frequentemente in una maggiore instabilità degli stati interni e in una ridotta capacità di recupero dell'equilibrio in seguito a stimoli stressanti.

            A livello comportamentale, tali difficoltà si manifestano in una maggiore reattività agli stimoli, nella ridotta capacità di auto-consolazione e in difficoltà nella regolazione attentiva, con tendenza alla distraibilità e alla ridotta persistenza sul compito \cite{montagna2016,clark2008}. Inoltre, si osserva una marcata sensibilità alle caratteristiche dell'ambiente: contesti poco strutturati o altamente stimolanti possono amplificare le difficoltà regolative, mentre ambienti prevedibili e relazioni di cura sensibili svolgono una funzione protettiva \cite{tronick2007}. In questo quadro, è importante sottolineare come la prematurità non conduca a esiti uniformi, ma si associ a traiettorie evolutive eterogenee, in cui fattori di rischio e di protezione contribuiscono a modulare l'organizzazione dei processi di regolazione nel tempo.

            In questa prospettiva, le difficoltà regolative nei bambini nati pretermine possono essere comprese come espressione di una vulnerabilità evolutiva che emerge dall'interazione tra fattori biologici precoci e qualità delle esperienze ambientali, influenzando le successive traiettorie di adattamento, come evidenziato da studi longitudinali e meta-analisi sugli esiti cognitivi e comportamentali a più lungo termine \cite{allotey2018}.

            Nel complesso, l'analisi della disregolazione in queste condizioni evidenzia come i processi di regolazione costituiscano un punto di osservazione privilegiato per comprendere le traiettorie di sviluppo nei primi anni di vita. Le difficoltà regolative emergono infatti dall'interazione tra caratteristiche neurobiologiche, risorse individuali e contesti di sviluppo, contribuendo a definire le modalità con cui il bambino organizza il proprio comportamento e si adatta alle richieste dell'ambiente nel corso dello sviluppo, e fornendo indicatori rilevanti per l'identificazione precoce delle traiettorie di rischio. In questo senso, tali difficoltà non si associano necessariamente a esiti disadattivi, ma possono configurarsi come indicatori di vulnerabilità la cui espressione dipende dall'interazione tra fattori di rischio e risorse di protezione nel corso dello sviluppo \cite{cicchetti2009}.

        \subsection{Disregolazione nei disturbi emotivi e comportamentali}
        \label{sec:2_3_3}

            All'interno della prospettiva transdiagnostica delineata nelle sezioni precedenti, i disturbi emotivi e comportamentali rappresentano un ambito particolarmente rilevante per approfondire le modalità attraverso cui la disregolazione si manifesta nei primi anni di vita. Tuttavia, a differenza di quanto osservato nei disturbi del neurosviluppo, in questo dominio emergono alcune specificità legate sia alla natura delle manifestazioni psicopatologiche sia alle difficoltà metodologiche connesse alla loro identificazione nella prima infanzia.

            In particolare, la letteratura evidenzia come nei bambini in età prescolare le manifestazioni di tipo esternalizzante risultino più facilmente osservabili e maggiormente indagate rispetto a quelle internalizzanti. Tale asimmetria non riflette necessariamente una reale differenza nella loro incidenza, ma è in parte riconducibile a limiti metodologici nella valutazione del funzionamento emotivo nei primi anni di vita. Le difficoltà internalizzanti, infatti, tendono a manifestarsi attraverso segnali meno osservabili e meno facilmente rilevabili mediante strumenti di assessment basati prevalentemente su osservazioni comportamentali o report esterni, risultando quindi più difficili da identificare in modo affidabile. Inoltre, la minore differenziazione fenomenologica che caratterizza la prima infanzia contribuisce a rendere meno chiari i confini tra espressioni normative e possibili segnali di difficoltà, rafforzando la necessità di una lettura dimensionale del funzionamento regolativo, coerente con la prospettiva transdiagnostica delineata nelle sezioni precedenti \cite{cicchetti2009,cole2004,astle2022}. Tale impostazione trova ulteriore supporto in evidenze empiriche che documentano la continuità tra funzionamento normativo e atipico nei primi anni di vita, evidenziando come le difficoltà regolative si distribuiscano lungo un continuum piuttosto che configurarsi in categorie discrete \cite{egger2006}.

            In questo contesto, i disturbi del comportamento e le manifestazioni esternalizzanti costituiscono il dominio maggiormente studiato nei primi anni di vita. Comportamenti quali oppositività, impulsività, difficoltà nel rispetto delle regole e nella modulazione della frustrazione si configurano come espressioni relativamente evidenti di difficoltà nella regolazione comportamentale ed emotiva. In questa prospettiva, tali manifestazioni possono essere interpretate come configurazioni di funzionamento in cui la disregolazione si esprime prevalentemente attraverso una difficoltà nel modulare risposte impulsive e stati emotivi ad alta attivazione, in linea con i modelli che sottolineano il ruolo dell'integrazione tra processi reattivi e sistemi di controllo \cite{beauchaine2019} .

            Evidenze empiriche indicano che difficoltà precoci nella regolazione emotiva e comportamentale risultano associate a una maggiore probabilità di sviluppare traiettorie esternalizzanti nel corso dell'età prescolare e scolare \cite{eisenberg2010,calkins2007hill,beauchaine2019}. In particolare, la disregolazione emotiva è stata identificata come un fattore centrale nello sviluppo dei problemi comportamentali, configurandosi come un meccanismo di rischio trasversale a diversi quadri psicopatologici \cite{beauchaine2015,graziano2016}.

            Tali manifestazioni risultano strettamente connesse alle componenti comportamentali ed esecutive del sistema regolativo, già descritte nel Capitolo 1 come dimensioni fondamentali dell'autoregolazione. In questo senso, la difficoltà nel controllo inibitorio, nella gestione dell'attivazione emotiva e nell'adattamento flessibile alle richieste del contesto può essere letta come espressione di una più ampia disregolazione che attraversa domini multipli del funzionamento.

            In questo quadro, le difficoltà osservabili nei disturbi del comportamento condividono importanti elementi di continuità con quanto descritto nei disturbi del neurosviluppo, e in particolare nel disturbo da deficit di attenzione/iperattività (ADHD), già approfondito nella sezione 2.3.1. In entrambi i casi, infatti, la disregolazione può essere interpretata come espressione di una difficoltà nel coordinamento tra processi di controllo e sistemi reattivi, che si manifesta attraverso impulsività, difficoltà nella modulazione dell'attenzione e una ridotta capacità di gestire stati emotivi ad alta attivazione. Tale continuità è supportata da evidenze empiriche che documentano come difficoltà nei processi di regolazione, in particolare a livello esecutivo ed emotivo, rappresentino un fattore condiviso tra diversi quadri del neurosviluppo e del comportamento \cite{sonugabarke2016,zelazo2020}.

            Tuttavia, una lettura esclusivamente centrata sulle manifestazioni esternalizzanti rischia di offrire una rappresentazione parziale del ruolo della disregolazione nei disturbi emotivi e comportamentali, come evidenziato dalla letteratura che sottolinea la natura multidimensionale dei processi regolativi \cite{beauchaine2015,cole2004}.

            In particolare, nei primi anni di vita risulta più complesso identificare e differenziare le manifestazioni internalizzanti, le quali tendono a presentarsi in forme meno osservabili e meno direttamente riconducibili a categorie diagnostiche definite \cite{egger2006}. In questa prospettiva, piuttosto che riferirsi a disturbi internalizzanti pienamente strutturati, appare più appropriato considerare la presenza di tratti precoci che possono configurarsi come potenziali precursori di successive traiettorie evolutive. Tra questi, un costrutto particolarmente rilevante è rappresentato dall'inibizione comportamentale, intesa come tendenza stabile a rispondere con ritiro, cautela e ipercontrollo in situazioni nuove o socialmente rilevanti \cite{kagan1987,fox2005}. Nei bambini piccoli, tali caratteristiche non configurano di per sé una condizione patologica, ma rappresentano una modalità di funzionamento che può essere associata a specifiche difficoltà nei processi di regolazione emotiva \cite{degnan2007}.

            A differenza delle manifestazioni esternalizzanti, in cui la disregolazione si esprime prevalentemente attraverso comportamenti impulsivi e difficoltà di controllo dell'azione, nei profili caratterizzati da inibizione comportamentale essa assume forme qualitativamente differenti. In questi casi, la difficoltà regolativa riguarda in misura maggiore la gestione degli stati emotivi interni, in particolare quelli connessi all'attivazione ansiosa e alla reattività sociale. Il bambino può mostrare una marcata difficoltà nel modulare l'arousal emotivo in contesti interpersonali, accompagnata da strategie di evitamento, rigidità comportamentale o ipercontrollo, che riflettono una modalità di regolazione meno flessibile e adattiva \cite{rubin2009}.

            Studi longitudinali mostrano come tratti precoci quali l'inibizione comportamentale siano associati a un aumentato rischio di sviluppare difficoltà internalizzanti, in particolare disturbi d'ansia \cite{kagan1987,fox2005,chronistuscano2009}. In questa prospettiva, l'inibizione comportamentale e il ritiro sociale sono stati ampiamente studiati come precursori di traiettorie internalizzanti, evidenziando una continuità tra caratteristiche temperamentali precoci e successivi esiti emotivi \cite{degnan2007,rubin2009}.

            Questa distinzione tra modalità di disregolazione disinibita e ipercontrollata risulta coerente con le prospettive dimensionali della psicopatologia, che individuano nelle dimensioni di internalizzazione ed esternalizzazione due principali organizzazioni del funzionamento psicologico, entrambe riconducibili a processi regolativi sottostanti \cite{zelazo2020}. In tale quadro, la disregolazione non può essere intesa unicamente come eccesso di espressività comportamentale, ma deve essere considerata come una difficoltà più ampia nella modulazione degli stati interni, che può manifestarsi sia in forme di discontrollo sia in modalità di eccessivo controllo.

            Dal punto di vista evolutivo, tali configurazioni assumono caratteristiche differenti nelle diverse fasi della prima infanzia. Nella fascia 0--3 anni, le manifestazioni regolative risultano ancora poco differenziate e fortemente dipendenti dal contesto e dai processi di co-regolazione, rendendo più difficile distinguere pattern specifici di funzionamento \cite{calkins2007hill,cole2004}. Con l'ingresso nell'età prescolare (3--6 anni), si osserva invece una progressiva maggiore stabilità e organizzazione dei comportamenti, che consente una più chiara identificazione delle diverse modalità di regolazione, incluse quelle riconducibili alle dimensioni internalizzanti ed esternalizzanti. In questo passaggio evolutivo, i tratti precoci, come l'inibizione comportamentale, possono assumere una maggiore coerenza e continuità nel tempo, configurandosi come indicatori più leggibili delle traiettorie di sviluppo \cite{cicchetti2002,sroufe2005}.

            In questa prospettiva, tali configurazioni possono essere interpretate come espressioni precoci di traiettorie evolutive divergenti, che, pur emergendo da una base comune di funzionamento regolativo, possono differenziarsi progressivamente nel corso dello sviluppo in relazione all'interazione tra caratteristiche individuali e contesto.

            Nel complesso, le evidenze empiriche disponibili supportano l'ipotesi che la disregolazione rappresenti un meccanismo condiviso che contribuisce a differenti esiti psicopatologici, configurandosi come una dimensione centrale nella prospettiva transdiagnostica \cite{beauchaine2019,astle2022}. Essa si manifesta infatti secondo modalità differenti, disinibite o ipercontrollate, ma riconducibili a una comune difficoltà nella modulazione dei processi emotivi, attentivi e comportamentali, rafforzando il suo ruolo nella comprensione dello sviluppo tipico e atipico nella prima infanzia.

        \subsection{Disregolazione come fattore di rischio transdiagnostico}
        \label{sec:2_3_4}

            Alla luce delle dimensioni e delle caratteristiche del funzionamento disregolato discusse nelle sezioni precedenti, la disregolazione nella prima infanzia può essere ulteriormente compresa non solo come espressione di un funzionamento regolativo meno organizzato, ma come una~dimensione di vulnerabilità a carattere transdiagnostico, rilevante per interpretare l'eterogeneità degli esiti evolutivi. In questa prospettiva, diversi contributi nell'ambito della psicopatologia dello sviluppo e dei modelli dimensionali sottolineano come specifici processi di base possano contribuire, in modo trasversale, all'emergere di differenti quadri psicopatologici \cite{astle2022,beauchaine2019}.

            Nel presente lavoro, per fattore di rischio transdiagnostico~si intende un processo o una dimensione funzionale che può contribuire all'emergere di differenti esiti psicopatologici, indipendentemente dalle specifiche categorie diagnostiche, configurandosi come una vulnerabilità condivisa tra più quadri clinici. Tale definizione si colloca all'interno di un approccio evolutivo che enfatizza la continuità tra sviluppo tipico e atipico e il ruolo dei processi regolativi nell'organizzazione del funzionamento nel tempo \cite{rutter1989,sroufe2009}.

            Come evidenziato nel Capitolo 1, i processi di regolazione nella prima infanzia si configurano come un sistema multidimensionale e progressivamente organizzato, che emerge dall'integrazione tra componenti emotive, attentive e comportamentali \cite{cole2004,eisenberg2010}. In questa cornice, difficoltà precoci nella modulazione di tali componenti possono essere intese come indicatori di un minor grado di integrazione del sistema regolativo, la cui organizzazione risulta ancora in fase di consolidamento \cite{cicchetti2009}. I diversi aspetti attraverso cui la disregolazione si manifesta, traiettorie evolutive divergenti, comorbidità, stabilità e continuità nel tempo, possono essere letti come espressioni interconnesse di questo funzionamento regolativo.

            Un primo elemento rilevante riguarda la relazione tra disregolazione e traiettorie evolutive divergenti. Evidenze provenienti da studi longitudinali nell'ambito della psicopatologia dello sviluppo suggeriscono che pattern precoci di difficoltà nella regolazione emotiva e comportamentale possano associarsi, nel tempo, a esiti differenti, anche molto eterogenei tra loro \cite{cicchetti2009,eisenberg2010,sroufe2009}. In questa prospettiva, la disregolazione può essere intesa non tanto come precursore specifico di un singolo disturbo, quanto come una~vulnerabilità generale, che può esprimersi attraverso configurazioni sintomatologiche diverse nel corso dello sviluppo. Ad esempio, difficoltà nella regolazione emotiva possono associarsi sia a manifestazioni internalizzanti, quali ansia o ritiro sociale, sia a problematiche esternalizzanti, come impulsività o comportamenti oppositivi \cite{beauchaine2019}. Tale variabilità può essere interpretata alla luce dell'interazione tra caratteristiche individuali e condizioni ambientali, che contribuiscono a modulare le traiettorie evolutive. Tali traiettorie possono essere ulteriormente comprese alla luce del ruolo svolto dai contesti di sviluppo, e in particolare dalla qualità delle interazioni precoci e delle opportunità di co-regolazione, che contribuiscono a modulare l'organizzazione del funzionamento regolativo nel tempo.

            Un secondo aspetto riguarda il fenomeno della comorbidità\textbf{,} frequentemente osservato nei disturbi del neurosviluppo. Studi epidemiologici e clinici evidenziano come la co-occorrenza di difficoltà appartenenti a domini diversi rappresenti una condizione comune, difficilmente riconducibile a modelli categoriali rigidi \cite{astle2022}. In questa prospettiva, la disregolazione può essere considerata come un possibile processo sottostante condiviso, che contribuisce alla sovrapposizione tra diversi quadri clinici. Le difficoltà nella modulazione delle emozioni, dell'attenzione e del comportamento implicano infatti il coinvolgimento di sistemi neurocognitivi interconnessi, la cui organizzazione può risultare meno stabile. In tale senso, modelli contemporanei sottolineano come le difficoltà di self-regulation rappresentino un nucleo centrale di vulnerabilità transdiagnostica, in grado di contribuire all'emergere di diversi quadri psicopatologici \cite{nigg2017}. In particolare, ricerche nell'ambito dello sviluppo neurocognitivo suggeriscono che difficoltà nelle funzioni esecutive, intese come processi di controllo top-down che supportano la regolazione del comportamento e delle emozioni e fungono da ponte tra componenti cognitive ed emotive, possano essere riscontrate trasversalmente in diversi quadri psicopatologici \cite{zelazo2020}. Allo stesso modo, la disregolazione emotiva è stata descritta come una dimensione comune a differenti domini della psicopatologia \cite{beauchaine2019}.

            Un ulteriore elemento riguarda la stabilità relativa e la capacità predittiva delle difficoltà regolative nel corso dello sviluppo\textbf{,} in relazione alla loro continuità con esiti successivi. Studi longitudinali e prospettici indicano che, sebbene la prima infanzia sia caratterizzata da una significativa variabilità nel funzionamento regolativo, pattern di disregolazione contraddistinti da elevata intensità, frequenza e pervasività tendono a mostrare una maggiore continuità nel tempo \cite{calkins2007,eisenberg2010}. Tale continuità non implica una traiettoria deterministica, ma può essere interpretata in termini di maggiore probabilità di esiti disadattivi, soprattutto in presenza di fattori di rischio contestuali \cite{rutter1989}.

            In particolare, evidenze provenienti da studi longitudinali suggeriscono che difficoltà nei processi di controllo inibitorio e nelle funzioni esecutive emergenti possano associarsi a esiti successivi meno adattivi sul piano comportamentale, emotivo e sociale \cite{zelazo2020}, evidenziando il ruolo del funzionamento regolativo nel sostenere l'adattamento nel tempo. In linea con tali evidenze, un ampio corpus di studi prospettici evidenzia una continuità tra difficoltà regolative nei primi anni di vita e manifestazioni psicopatologiche nelle fasi successive dello sviluppo, sia in termini omotipici sia eterotipici \cite{beauchaine2019,eisenberg2010,rutter1989}.

            Ad esempio, una disregolazione emotiva precoce può associarsi, nel tempo, a sintomi ansiosi o depressivi, mentre difficoltà nella regolazione comportamentale possono essere osservate in associazione a problematiche esternalizzanti o a difficoltà nel controllo degli impulsi. Allo stesso tempo, uno stesso pattern iniziale può essere associato a esiti differenti, in funzione dell'interazione tra fattori individuali e contestuali.

            Nel complesso, tali evidenze supportano una concettualizzazione della disregolazione come~dimensione trasversale del funzionamento, che può contribuire alla configurazione di differenti traiettorie evolutive e psicopatologiche nel corso dello sviluppo \cite{beauchaine2019,rutter1989,nigg2017}. In questa prospettiva, il funzionamento regolativo può essere inteso come un elemento centrale nell'organizzazione dello sviluppo, il cui grado di integrazione risulta rilevante per comprendere la continuità tra difficoltà precoci ed esiti successivi. Alla luce di ciò, appare utile approfondire i fattori che modulano tali traiettorie e i processi attraverso cui il sistema regolativo si organizza nel tempo, come verrà discusso nella sezione successiva.

    \section{Modelli teorici del funzionamento transdiagnostico}
    \label{sec:2_4}

        Alla luce delle evidenze teoriche ed empiriche discusse nelle sezioni precedenti, emerge la necessità di una cornice teorica capace di spiegare la natura trasversale delle difficoltà osservate nello sviluppo. In questo contesto, il crescente interesse per una lettura transdiagnostica del funzionamento psicologico si colloca all'interno di un più ampio processo di revisione dei modelli tradizionali della psicopatologia dello sviluppo \cite{nolenhoeksema2011,astle2022}. In particolare, la difficoltà di ricondurre la complessità e l'eterogeneità delle manifestazioni cliniche entro categorie diagnostiche discrete ha progressivamente orientato la ricerca verso l'identificazione di dimensioni di funzionamento e di meccanismi di base condivisi tra differenti traiettorie evolutive.

        In questo quadro, un aspetto centrale riguarda l'individuazione di processi fondamentali che contribuiscono all'organizzazione del funzionamento psicologico lungo lo sviluppo. Tra questi, i meccanismi di autoregolazione assumono un ruolo particolarmente rilevante, in quanto implicati nella modulazione delle emozioni, dell'attenzione e del comportamento e potenzialmente coinvolti in differenti configurazioni psicopatologiche, in linea con modelli dello sviluppo che sottolineano il ruolo dei processi di controllo esecutivo e regolazione \cite{diamond2013,zelazo2015}.

        Come discusso nel Capitolo 1, l'autoregolazione può essere concepita come un sistema complesso e multilivello, derivante dall'integrazione tra processi emotivi, cognitivi, fisiologici e comportamentali. In tale prospettiva, la possibilità di considerarla come una dimensione rilevante in chiave transdiagnostica si fonda sulla sua implicazione nei processi di adattamento del comportamento lungo il continuum tra sviluppo tipico e atipico.

        Parallelamente, i sistemi diagnostici categoriali, pur mantenendo una rilevanza in ambito clinico e descrittivo, risultano limitati nel cogliere la natura dinamica, multidimensionale e contestuale dei processi di sviluppo \cite{kotov2017}, soprattutto nelle fasi precoci della vita. Tale limite ha favorito uno spostamento dell'attenzione dall'identificazione di configurazioni sintomatologiche discrete all'analisi dei processi sottostanti che contribuiscono all'organizzazione del funzionamento psicologico \cite{nolenhoeksema2011}.

        All'interno di questo cambiamento teorico, i modelli contemporanei hanno progressivamente abbandonato una concezione rigidamente categoriale della psicopatologia, orientandosi verso approcci dimensionali e gerarchici che consentono di descrivere il funzionamento psicologico lungo continuum di variabilità. In particolare, tali prospettive hanno trovato espressione nei modelli gerarchici della psicopatologia, tra cui i modelli bifattoriali che includono un fattore generale di psicopatologia (p factor), inteso come indice di vulnerabilità condivisa tra differenti disturbi \cite{caspi2014}, accanto a dimensioni più specifiche, come quelle internalizzanti ed esternalizzanti \cite{kotov2017}.

        Tali modelli riflettono un cambiamento più profondo nel modo di concepire i meccanismi sottostanti lo sviluppo tipico e atipico, ponendo le basi per l'identificazione di dimensioni trasversali e processi condivisi tra differenti condizioni cliniche. In questa prospettiva, l'approccio transdiagnostico si configura come un tentativo di superare i limiti delle classificazioni tradizionali, focalizzandosi su caratteristiche e meccanismi che attraversano i confini diagnostici e che risultano maggiormente rappresentativi della complessità del funzionamento individuale \cite{nolenhoeksema2011,astle2022}.

        In questo quadro teorico, l'analisi dei modelli del funzionamento transdiagnostico consente di comprendere come le difficoltà osservabili nei diversi ambiti della psicopatologia possano essere ricondotte, almeno in parte, a disfunzioni nei sistemi di base che sostengono la regolazione dell'esperienza e del comportamento. Tra questi, i processi di autoregolazione emergono come un nodo centrale, contribuendo all'organizzazione delle risposte emotive, attentive e comportamentali in relazione alle richieste del contesto e risultando implicati in molteplici traiettorie evolutive.

        Alla luce di queste considerazioni, risulta necessario approfondire i limiti dei modelli categoriali tradizionali e comprendere come tali criticità abbiano favorito l'emergere di approcci dimensionali e transdiagnostici. In questa prospettiva, l'analisi dei diversi modelli teorici consentirà di chiarire in che modo i processi di autoregolazione possano essere concepiti come una dimensione trasversale rilevante per la comprensione delle traiettorie di sviluppo tipico e atipico.

        \subsection{Limiti dei modelli categoriali nella psicopatologia dello sviluppo}
        \label{sec:2_4_1}

            I modelli categoriali della psicopatologia, formalizzati nei principali sistemi diagnostici (come il DSM-5), si basano sull'assunto che i disturbi possano essere identificati come entità discrete, caratterizzate da configurazioni sintomatologiche relativamente specifiche e distinguibili tra loro. Sebbene tale approccio abbia fornito un quadro di riferimento utile per la classificazione e la comunicazione clinica, numerose evidenze empiriche ne hanno progressivamente messo in luce alcuni limiti significativi, in particolare nell'ambito della psicopatologia dello sviluppo.

            Tra le principali criticità emerge l'elevata eterogeneità osservabile all'interno delle singole categorie diagnostiche. Bambini che condividono la stessa diagnosi possono infatti presentare profili di funzionamento molto diversi in termini di caratteristiche cognitive, emotive e comportamentali, rendendo difficile identificare pattern coerenti e stabili che definiscano in modo univoco il disturbo. Tale variabilità riflette la natura complessa e multidimensionale dei processi di sviluppo, come evidenziato nella letteratura sulle traiettorie evolutive \cite{rutter1989}.

            A questo si affianca la frequente co-occorrenza tra disturbi differenti. Le difficoltà evolutive tendono infatti a presentarsi in forma sovrapposta piuttosto che isolata, con elevati livelli di comorbilità tra condizioni diverse. Questo dato mette in discussione l'idea che i disturbi rappresentino entità indipendenti e suggerisce piuttosto l'esistenza di processi che attraversano differenti quadri clinici \cite{astle2022}.

            Un ulteriore limite riguarda la difficoltà dei modelli categoriali nel rappresentare la continuità tra sviluppo tipico e atipico. Come discusso nelle sezioni precedenti, il funzionamento psicologico nei primi anni di vita si colloca lungo un continuum di variabilità, in cui le differenze riguardano il grado di organizzazione e adattività dei processi regolativi, piuttosto che la presenza o assenza di specifiche categorie diagnostiche. In una fase evolutiva caratterizzata da competenze ancora emergenti e da una scarsa differenziazione delle manifestazioni comportamentali, l'applicazione di criteri categoriali rigidi risulta pertanto particolarmente problematica.

            Considerati nel loro insieme, questi elementi mettono in evidenza una criticità più generale dei modelli categoriali, ovvero la loro difficoltà nel cogliere i processi di base che sottendono il funzionamento psicologico. La focalizzazione su insiemi di sintomi osservabili tende, infatti, a trascurare i meccanismi attraverso cui tali manifestazioni si organizzano e si modulano nel tempo. In particolare, tali modelli risultano strutturalmente limitati nel rappresentare il ruolo dei sistemi di regolazione, i quali costituiscono un livello esplicativo centrale per la comprensione dell'organizzazione del funzionamento psicologico nello sviluppo e delle traiettorie che conducono a esiti tipici e atipici \cite{beauchaine2019}.

            Nel complesso, questi limiti hanno contribuito a mettere in discussione l'adeguatezza dei modelli categoriali nel descrivere la complessità del funzionamento psicologico nello sviluppo, favorendo l'emergere di approcci alternativi orientati all'analisi dimensionale e all'identificazione di processi transdiagnostici. Questi aspetti verranno approfonditi nella sezione seguente.

        \subsection{Approcci dimensionali emergenti nello studio della psicopatologia}
        \label{sec:2_4_2}

            In risposta ai limiti dei modelli categoriali descritti nella sezione precedente, la ricerca in psicopatologia dello sviluppo si è progressivamente orientata verso approcci dimensionali, che consentono di descrivere il funzionamento psicologico lungo un continuum di variabilità piuttosto che attraverso categorie discrete. In tale prospettiva, le manifestazioni psicopatologiche possono essere comprese non come entità qualitativamente distinte, ma come espressioni differenti, per grado e organizzazione, di dimensioni sottostanti che caratterizzano il funzionamento individuale.

            In questo senso, le dimensioni possono essere intese come assi di variazione lungo i quali gli individui differiscono per livello di funzionamento, riflettendo l'organizzazione di processi psicologici sottostanti piuttosto che la presenza di categorie discrete.

            L'analisi delle dimensioni non si limita quindi a una funzione descrittiva, ma si configura progressivamente come uno strumento per comprendere i processi di base che organizzano il funzionamento psicologico lungo lo sviluppo, in linea con prospettive che sottolineano il ruolo dei meccanismi transdiagnostici nella spiegazione della psicopatologia \cite{nolenhoeksema2011}.

            Un primo elemento distintivo degli approcci dimensionali riguarda l'assunzione di una continuità tra sviluppo tipico e atipico. Le differenze osservabili nei comportamenti emotivi, cognitivi e comportamentali possono essere interpretate in termini di variazioni lungo dimensioni condivise, piuttosto che come esiti dicotomici. Questa prospettiva consente di cogliere in modo più adeguato la variabilità individuale e la natura dinamica dei processi di sviluppo, in particolare nella prima infanzia, fase caratterizzata da elevata plasticità e da una progressiva organizzazione dei sistemi di regolazione, che si manifestano con ampia variabilità in relazione alle condizioni contestuali \cite{zelazo2020,beauchaine2019}.

            In questo quadro, l'attenzione tende a spostarsi dalla descrizione delle configurazioni sintomatologiche all'analisi dei processi sottostanti che contribuiscono alla strutturazione del funzionamento psicologico. Le dimensioni considerate rilevanti non si limitano quindi a rappresentare aggregati di sintomi, ma possono essere intese come espressione di modalità di funzionamento più profonde, connesse alla capacità dell'individuo di modulare, coordinare e integrare diversi sistemi, tra cui quelli emotivi, attentivi e comportamentali, in linea con prospettive che evidenziano il ruolo centrale dei sistemi di regolazione nello sviluppo e nella psicopatologia \cite{zelazo2020,beauchaine2019}. In questa prospettiva, tali dimensioni riflettono processi di funzionamento che contribuiscono alla strutturazione del comportamento lungo lo sviluppo e alla variabilità delle traiettorie evolutive, manifestandosi trasversalmente a differenti configurazioni psicopatologiche.

            Questo cambiamento di prospettiva trova espressione nei modelli dimensionali e gerarchici della psicopatologia, che organizzano le diverse manifestazioni cliniche all'interno di strutture multilivello. In particolare, i modelli bifattoriali ipotizzano la presenza di un fattore generale di psicopatologia (p factor), inteso come indice di vulnerabilità condivisa tra differenti condizioni cliniche \cite{caspi2014}, accanto a dimensioni più specifiche che riflettono pattern di funzionamento più circoscritti. Analogamente, i modelli gerarchici descrivono l'organizzazione dei sintomi lungo livelli progressivi di generalità, nei quali dimensioni più ampie si articolano in sottodimensioni più specifiche, consentendo di rappresentare sia la comunanza sia la specificità delle diverse manifestazioni psicopatologiche \cite{kotov2017}.

            Secondo prospettive che promuovono un superamento dei modelli categoriali tradizionali \cite{nolenhoeksema2011}, tra i modelli che hanno contribuito a sostenere una lettura dimensionale della psicopatologia si colloca il Research Domain Criteria (RDoC), proposto dal National Institute of Mental Health, che analizza i disturbi mentali attraverso dimensioni di funzionamento neurobiologico e comportamentale \cite{cuthbert2013}. Tuttavia, tale approccio, prevalentemente orientato a livelli neurobiologici, non è stato sviluppato specificamente per lo studio delle prime fasi dello sviluppo e richiede pertanto un'integrazione con modelli propri della psicopatologia dello sviluppo.

            Nel complesso, il passaggio da modelli categoriali ad approcci dimensionali riflette un cambiamento più ampio nel modo di concepire la psicopatologia, orientato all'identificazione di processi di base che contribuiscono all'organizzazione del funzionamento psicologico. Se tali modelli consentono di descrivere la struttura della variabilità psicopatologica, essi aprono al contempo alla necessità di comprenderne i meccanismi sottostanti. In questo senso, i modelli transdiagnostici rappresentano un ulteriore sviluppo teorico, volto a identificare i processi che attraversano differenti traiettorie evolutive e configurazioni psicopatologiche.

        \subsection{Modelli transdiagnostici nello sviluppo precoce: vulnerabilità condivise e integrazione dei sistemi}
        \label{sec:2_4_3}

            Se gli approcci dimensionali consentono di descrivere la~struttura della variabilità del funzionamento psicologico, i modelli transdiagnostici mirano a spiegare i meccanismi che la generano\textbf{,} individuando i processi che contribuiscono in modo trasversale allo sviluppo di differenti esiti psicopatologici. In continuità con i modelli gerarchici della psicopatologia, che includono anche il fattore generale di psicopatologia come indice di varianza condivisa tra differenti manifestazioni cliniche \cite{caspi2014,kotov2017}, tale variabilità può essere interpretata non tanto come una dimensione astratta, quanto come espressione di processi funzionali che contribuiscono in modo coordinato alla strutturazione del funzionamento psicologico \cite{nolenhoeksema2011,dalgleish2020}. In questa prospettiva, la psicopatologia è concepita come esito di sistemi di interazione dinamica tra sintomi e processi sottostanti, piuttosto che come manifestazione di entità latenti discrete.

            Questo orientamento trova una formulazione esplicita nei modelli transdiagnostici process-based, che interpretano le manifestazioni psicopatologiche come esiti di alterazioni in funzioni psicologiche di base comuni a differenti condizioni cliniche \cite{nolenhoeksema2011,dalgleish2020}. In tale quadro, la variabilità osservabile nei profili di funzionamento viene ricondotta a configurazioni differenti di vulnerabilità che coinvolgono sistemi implicati nella regolazione del comportamento, piuttosto che a categorie diagnostiche discrete.

            Un contributo particolarmente rilevante in questa direzione proviene dai modelli transdiagnostici della regolazione emotiva, che individuano nella disregolazione una dimensione centrale implicata nello sviluppo di molteplici condizioni psicopatologiche \cite{beauchaine2015,beauchaine2019}. In tali prospettive, la difficoltà nel modulare l'intensità, la durata e l'espressione delle emozioni non rappresenta un fenomeno specifico di singoli disturbi, ma una dimensione trasversale che, in interazione con altri fattori, può contribuire alla strutturazione di esiti eterogenei lungo lo sviluppo.

            In questa stessa direzione, i modelli neuroevolutivi delle funzioni esecutive evidenziano come i sistemi di controllo cognitivo costituiscano componenti fondamentali per la regolazione del comportamento e dell'esperienza, risultando implicati trasversalmente in un'ampia gamma di difficoltà evolutive \cite{zelazo2015,zelazo2020,diamond2013}. Il funzionamento adattivo dipende infatti dalla capacità di integrare dinamiche top-down, sostenute dalle funzioni esecutive, con dinamiche bottom-up legate alla reattività emotiva e fisiologica, evidenziando il carattere multilivello dei sistemi regolativi.

            Parallelamente, modelli transdiagnostici dimensionali recenti, sviluppati nell'ambito dei disturbi del neurosviluppo, sottolineano come le difficoltà osservabili nei diversi quadri clinici riflettano combinazioni eterogenee di dimensioni cognitive e comportamentali, che tendono a organizzarsi in profili funzionali indipendenti dai confini diagnostici \cite{astle2022}. In queste prospettive, domini quali l'attenzione, il controllo esecutivo e la regolazione emotiva vengono concepiti come dimensioni fondamentali la cui variabilità contribuisce alla definizione di differenti configurazioni di sviluppo.

            Considerati congiuntamente, questi modelli convergono nell'individuare alcune aree di vulnerabilità che attraversano le diverse traiettorie psicopatologiche, in particolare difficoltà nel controllo inibitorio, nella regolazione emotiva e nella gestione dell'attenzione. Tuttavia, essi differiscono per il livello di analisi privilegiato e per il peso attribuito ai diversi sistemi coinvolti, evidenziando la necessità di un'integrazione tra prospettive diverse per comprendere i meccanismi transdiagnostici. Un limite condiviso da queste prospettive risiede inoltre nella difficoltà di descrivere in modo unitario le modalità attraverso cui tali processi interagiscono dinamicamente, soprattutto nelle fasi precoci dello sviluppo.

            La rilevanza di tali vulnerabilità risulta più chiara se considerata alla luce dell'integrazione tra i diversi sistemi che sostengono il funzionamento regolativo. Come discusso nel Capitolo 1, l'autoregolazione emerge dall'interazione tra componenti emotive, cognitive, comportamentali e fisiologiche, configurandosi come un sistema complesso e multilivello \cite{beauchaine2019,zelazo2015}. In questa prospettiva, il funzionamento regolativo può essere concepito come un sistema emergente, la cui organizzazione dipende dalla qualità del coordinamento tra i diversi domini, piuttosto che dal livello di competenza in ciascun dominio considerato isolatamente.

            I modelli transdiagnostici consentono quindi di interpretare le difficoltà evolutive come espressione di una alterazione nelle dinamiche di integrazione tra tali sistemi. In particolare, uno squilibrio tra sistemi di reattività e sistemi di controllo può compromettere la capacità del bambino di modulare in modo flessibile emozioni, attenzione e comportamento, contribuendo alla variabilità delle traiettorie evolutive \cite{beauchaine2019,zelazo2015}.

            In questo quadro teorico, l'autoregolazione può essere intesa come il processo attraverso cui tali sistemi vengono coordinati e integrati, contribuendo all'organizzazione del comportamento in relazione alle richieste del contesto \cite{beauchaine2019}. In quanto sistema multidimensionale e integrato, essa risulta implicata trasversalmente nei diversi domini del funzionamento e contribuisce alla strutturazione del comportamento lungo differenti traiettorie evolutive.

            In questa prospettiva, l'autoregolazione può essere interpretata come un meccanismo di base che contribuisce all'organizzazione di molteplici traiettorie psicopatologiche, configurandosi come una dimensione transdiagnostica centrale nello sviluppo precoce \cite{dalgleish2020}.

            L'analisi dei modelli teorici del funzionamento transdiagnostico evidenzia quindi come l'autoregolazione rappresenti uno dei principali processi attraverso cui si articolano le dimensioni di vulnerabilità nello sviluppo. La sua natura multidimensionale e trasversale ne sottolinea la rilevanza per la comprensione delle traiettorie evolutive, evidenziando al contempo la complessità del suo inquadramento teorico ed empirico.

            Nel complesso, le evidenze discusse nel presente capitolo delineano l'autoregolazione come una dimensione centrale per la comprensione dello sviluppo nella prima infanzia, la cui variabilità si colloca lungo un continuum che va da modalità di funzionamento più adattive a configurazioni progressivamente meno organizzate \cite{calkins2007,cicchetti2009}. In questa prospettiva, la disregolazione può essere intesa come una dimensione trasversale che attraversa diversi domini del funzionamento e contribuisce alla variabilità delle traiettorie evolutive, configurandosi come possibile tratto a carattere transdiagnostico \cite{beauchaine2019,astle2021}.

            Tale interpretazione risulta coerente con le prospettive dimensionali della psicopatologia dello sviluppo, che pongono l'attenzione sui processi di base piuttosto che sulle categorie diagnostiche, in particolare nelle fasi precoci della vita \cite{caspi2014,kotov2017}. Tuttavia, la natura multidimensionale, dinamica e in parte latente dei processi regolativi rende complessa la loro osservazione e operazionalizzazione empirica, soprattutto nella finestra evolutiva tra i 18 e i 36 mesi \cite{cole2004,eisenberg2010}.

            In questa prospettiva, risulta quindi rilevante interrogarsi sulle modalità attraverso cui tali processi possano essere valutati in modo sistematico e sensibile alle caratteristiche dello sviluppo precoce. Questo aspetto verrà approfondito nel capitolo successivo, dedicato alla valutazione dell'autoregolazione nella prima infanzia.
